\chapter*{参考資料}
\markboth{参考資料}{}
\addcontentsline{toc}{chapter}{参考資料}
本書の説明と問題は,この構成に合わせて作成した.以下は,扱う範囲や標準的な定理の条件,公式の表記を確認した資料である.発展内容がすべて高校の学習範囲であることを意味するものではない.
\begin{enumerate}[leftmargin=*]
\item 文部科学省『高等学校学習指導要領（平成30年告示）解説 数学編・理数編』,2018年.数学II・IIIの極限・微分・積分の範囲の確認.掲載ページ：\url{https://www.mext.go.jp/a_menu/shotou/new-cs/1407074.htm}.
\item Gilbert Strang and Edwin Herman, \textit{Calculus Volume 1}, OpenStax, 2016. 第5章第3節「The Fundamental Theorem of Calculus」.基本定理の条件と面積関数.\url{https://openstax.org/books/calculus-volume-1/pages/5-3-the-fundamental-theorem-of-calculus}.
\item Gilbert Strang and Edwin Herman, \textit{Calculus Volume 2}, OpenStax, 2016. 第3章第3節「Trigonometric Substitution」.三角置換の範囲と符号.\url{https://openstax.org/books/calculus-volume-2/pages/3-3-trigonometric-substitution}.
\item NIST, \textit{Digital Library of Mathematical Functions}, \S5.12, ``Beta Function''. ベータ積分と三角関数の積分による表示.\url{https://dlmf.nist.gov/5.12}.

\item Galileo Galilei, \textit{Dialogues Concerning Two New Sciences}, 1638. Henry Crew and Alfonso de Salvio訳,1914. 第3日の一様な加速運動. バージニア大学公開の翻訳本文：\url{https://galileo.phys.virginia.edu/classes/109N/tns_draft/tns_160to243.html}.
\item Archimedes, \textit{Quadrature of the Parabola}, T. L. Heath訳,\textit{The Works of Archimedes},1897,命題18〜24. 三角形の面積の比較と取り尽くし法. 原著の英訳：\url{https://www.aproged.pt/biblioteca/worksofarchimede.pdf}.
\item Isaac Newton, ``History of the Method of Fluxions'', Cambridge University Library,MS Add.3968. 変化する量とその速度という説明を確認. The Newton Projectの翻刻：\url{https://www.newtonproject.ox.ac.uk/view/texts/diplomatic/NATP00356}.
\end{enumerate}
オンライン資料の確認日：2026年10月22日.本書で証明を省略して採用する事実は,該当箇所でその旨を示した.
